\documentclass[12pt]{article}
\usepackage[utf8]{inputenc}
\usepackage[T1]{fontenc}
\usepackage[margin=1in]{geometry}
\usepackage{setspace}
\usepackage{xcolor}
\usepackage[colorlinks=true, urlcolor=blue!60!black, linkcolor=blue!60!black]{hyperref}
\setlength{\parindent}{2em}
\setlength{\parskip}{1.25em}
\renewcommand{\baselinestretch}{1.0}
\title{\textit{United States v. Curtiss-Wright Export Corp.}, 299 U.S. 304 (1936)}
\author{Case Brief}
\date{}

\begin{document}
\maketitle

\section{Case Metadata}
\begin{description}
  \item[Citation:] 299 U.S. 304 (1936) (No.\ 98).
  \item[Argued / Decided:] Argued November 19--20, 1936; Decided December 21, 1936.
  \item[Vote:] 7--1 (Justice Harlan Fiske Stone took no part in the consideration or decision of the case).
  \item[Author:] Justice George Sutherland (joined by Chief Justice Hughes and Justices Van Devanter, Brandeis, Butler, Roberts, and Cardozo).
  \item[Concurrences / Dissents:] Justice James Clark McReynolds noted his dissent without opinion, stating that the District Court reached the right conclusion and its judgment ought to be affirmed. No concurring opinions.
  \item[Chapter / Topic:] Separation of Powers and the Executive (Foreign Affairs Power; Nondelegation in External Affairs; Executive Primacy).
  \item[Primary Source URLs:] \href{https://supreme.justia.com/cases/federal/us/299/304/}{Justia U.S. Supreme Court Center} (\url{https://supreme.justia.com/cases/federal/us/299/304/}); \href{https://www.law.cornell.edu/supremecourt/text/299/304}{Cornell LII} (\url{https://www.law.cornell.edu/supremecourt/text/299/304}).
\end{description}

\section{Facts and Procedural History}
During the Chaco War between Bolivia and Paraguay over the disputed Gran Chaco region, Congress enacted a Joint Resolution on May 28, 1934 (c.\ 365, 48 Stat.\ 811). Section~1 provided that if the President found that prohibiting the sale of arms and munitions of war in the United States to the countries engaged in armed conflict in the Chaco ``may contribute to the reestablishment of peace between those countries,'' and if---after consultation with other American Republics and with their cooperation, as well as that of other governments he deemed necessary---he issued a proclamation to that effect, it would become unlawful to sell arms or munitions of war in the United States to those belligerents, ``except under such limitations and exceptions as the President prescribes,'' until otherwise ordered by the President or Congress. Section~2 made violations a crime punishable by a fine of up to \$10,000, imprisonment for up to two years, or both.

That same day (May 28, 1934), President Franklin D.\ Roosevelt issued a proclamation (48 Stat.\ 1744) reciting the required findings, activating the arms embargo against Bolivia and Paraguay, and delegating to the Secretary of State the authority to prescribe exceptions and limitations. After hostilities subsided, President Roosevelt issued a second proclamation on November 14, 1935 (49 Stat.\ 3480), revoking the embargo effective November 29, 1935, subject to a proviso that revocation would not extinguish criminal liabilities incurred while the first proclamation was in force.

On January 27, 1936, a federal grand jury in the Southern District of New York indicted Curtiss-Wright Export Corp.\ and co-defendants for conspiring, beginning May 29, 1934, to sell 15 machine guns in the United States to Bolivia in violation of the Joint Resolution and the May 1934 proclamation. The defendants demurred to the indictment on three grounds: (1) the Joint Resolution unconstitutionally delegated legislative power to the President (invoking \textit{Panama Refining} and \textit{Schechter Poultry}); (2) the President failed to make essential jurisdictional findings in his proclamation; and (3) the November 1935 revocation extinguished any criminal liability because the Joint Resolution lacked a statutory saving clause.

The District Court (14 F.\ Supp.\ 230 (S.D.N.Y.\ 1936)) sustained the demurrer on the first ground---holding the Joint Resolution an invalid delegation of legislative power---and quashed the indictment, while overruling the second and third grounds. The United States took a direct appeal to the Supreme Court under the Criminal Appeals Act of 1907 (18 U.S.C.\ \S\,682).

\section{Issues Presented}
\begin{enumerate}
  \item \textbf{Foreign Affairs Nondelegation:} Did the May 28, 1934 Joint Resolution unconstitutionally delegate legislative power to the President by authorizing him to activate, limit, and terminate a criminal arms embargo on foreign belligerents?
  \item \textbf{Jurisdictional Findings \& Effect of Revocation:} Were the President's factual recitals in the May 1934 proclamation sufficient under the Joint Resolution, and did his November 1935 revocation extinguish criminal liability for violations committed while the embargo was in force?
\end{enumerate}

\section{Holding}
The Supreme Court reversed (7--1) and remanded for further proceedings. First, the Joint Resolution was not an unconstitutional delegation of legislative power because Congress may confer far broader discretion upon the President in foreign and external affairs than would be permissible in purely domestic matters. Second, the President's proclamation tracked the statutory language and satisfied every jurisdictional requirement. Third, revoking the first proclamation did not repeal the Joint Resolution or extinguish criminal liability for offenses committed during the embargo's operative period.

\section{Reasoning}
Writing for the Court, Justice Sutherland expressly assumed without deciding that the Joint Resolution's delegation might have been invalid if confined to internal domestic affairs. The decisive distinction, however, was that the Resolution's ``whole aim'' was ``to affect a situation entirely external to the United States and falling within the category of foreign affairs'' (299 U.S.\ at 315). Sutherland developed three interlocking rationales:

\begin{enumerate}
  \item \textbf{Extraconstitutional Origin of External Sovereignty:} Sutherland asserted a sharp dichotomy between internal and external federal powers. The rule of enumerated powers---that the federal government possesses only powers expressly granted by the Constitution or necessarily implied therefrom---is ``categorically true only in respect of our internal affairs,'' where the Constitution carved federal authority out of the preexisting powers of the States. Because the individual States ``severally never possessed international powers,'' external sovereignty passed directly from the British Crown upon independence to the Colonies in their collective corporate capacity as the United States of America. Consequently, federal authority over external affairs ``did not depend upon the affirmative grants of the Constitution,'' but vested as ``necessary concomitants of nationality'' (299 U.S.\ at 315--318).

  \item \textbf{The President as ``Sole Organ'' and Functional Institutional Competence:} Turning from federal power generally to executive power specifically, Sutherland argued that participation in external affairs is uniquely concentrated in the Presidency. Quoting Representative John Marshall's March 7, 1800 House floor speech, Sutherland declared that the President is ``the sole organ of the nation in its external relations, and its sole representative with foreign nations'' (299 U.S.\ at 319). Crucially, Sutherland emphasized that the Court was reviewing not merely delegated legislative authority, but statutory authority \textit{combined with} ``the very delicate, plenary and exclusive power of the President as the sole organ of the federal government in the field of international relations---a power which does not require as a basis for its exercise an act of Congress'' (though still subordinate to constitutional limits) (299 U.S.\ at 319--320).

  Functionally, Sutherland stressed that the President---not Congress---possesses confidential diplomatic and consular sources of information, unity of design, and the secrecy and dispatch necessary to avoid ``serious embarrassment'' in international relations. Requiring Congress to prescribe narrow, domestic-style statutory standards in foreign affairs would therefore be unwise and impracticable.

  \item \textbf{Unbroken Historical Practice:} Sutherland reinforced these structural arguments with an extensive catalog of congressional enactments dating back to 1794---including President Washington's 1794 embargo authority, early Non-Intercourse Acts, retaliatory tariff and tonnage statutes, and the 1898, 1912, and 1922 arms-embargo resolutions. This ``uniform, long-continued and undisputed legislative practice'' spanning a century and a half constituted a practical constitutional exposition entitled to decisive weight (299 U.S.\ at 322--329).

  \item \textbf{Jurisdictional Findings and Non-Abatement:} Exercising appellate jurisdiction to reach the alternative demurrer grounds, the Court held that the President's recital that he had consulted other American Republics and secured necessary cooperation tracked the Resolution and was sufficient. Finally, the November 1935 proclamation merely ended the factual condition that triggered the embargo prospectively; it did not repeal the underlying statute or abate prosecutions for crimes committed between May 28, 1934, and November 14, 1935.
\end{enumerate}

\section{Concurring and Dissenting Opinions}
\textbf{Justice McReynolds, dissenting:} Justice McReynolds issued a two-sentence dissent without a written opinion, stating simply that he disagreed and believed the District Court had reached the right conclusion in holding the delegation unconstitutional and ought to be affirmed (299 U.S.\ at 333). (Justice Stone did not participate.)

\section{Key Quotations}
\begin{quote}
``The broad statement that the federal government can exercise no powers except those specifically enumerated in the Constitution, and such implied powers as are necessary and proper to carry into effect the enumerated powers, is categorically true only in respect of our internal affairs.''\\
--- Justice Sutherland, 299 U.S.\ at 315--316.
\end{quote}

\begin{quote}
``As a result of the separation from Great Britain by the colonies, acting as a unit, the powers of external sovereignty passed from the Crown not to the colonies severally, but to the colonies in their collective and corporate capacity as the United States of America. \dots\ It results that the investment of the federal government with the powers of external sovereignty did not depend upon the affirmative grants of the Constitution.''\\
--- Justice Sutherland, 299 U.S.\ at 316, 318.
\end{quote}

\begin{quote}
``In this vast external realm, with its important, complicated, delicate and manifold problems, the President alone has the power to speak or listen as a representative of the nation. He makes treaties with the advice and consent of the Senate; but he alone negotiates.''\\
--- Justice Sutherland, 299 U.S.\ at 319.
\end{quote}

\begin{quote}
``It is important to bear in mind that we are here dealing not alone with an authority vested in the President by an exertion of legislative power, but with such an authority plus the very delicate, plenary and exclusive power of the President as the sole organ of the federal government in the field of international relations---a power which does not require as a basis for its exercise an act of Congress but which, of course, like every other governmental power, must be exercised in subordination to the applicable provisions of the Constitution.''\\
--- Justice Sutherland, 299 U.S.\ at 319--320.
\end{quote}

\begin{quote}
``[C]ongressional legislation which is to be made effective through negotiation and inquiry within the international field must often accord to the President a degree of discretion and freedom from statutory restriction which would not be admissible were domestic affairs alone involved.''\\
--- Justice Sutherland, 299 U.S.\ at 320.
\end{quote}

\section{Doctrinal Significance and Modern Legacy (Through 2026)}
\textit{Curtiss-Wright} has served for nine decades as the Executive Branch's premier citation for broad presidential primacy in foreign affairs---yet the Supreme Court has repeatedly trimmed its sweeping dicta while preserving its core holding on foreign-affairs delegations:

\begin{itemize}
  \item \textbf{\textit{Youngstown Sheet \& Tube Co. v. Sawyer}, 343 U.S.\ 579 (1952):} In his canonical tripartite concurrence, Justice Jackson reined in \textit{Curtiss-Wright} by classifying it as a \textbf{Zone~1} case (where the President acts with express congressional authorization) and dismissing much of Sutherland's expansive language as dicta. Jackson emphasized that \textit{Curtiss-Wright} involved ``the right of the President to act under and in accord with an Act of Congress,'' not a license to defy Congress or exercise unchecked domestic power under a foreign-affairs label (343 U.S.\ at 635--636 n.\,2).
  \item \textbf{\textit{Dames \& Moore v. Regan}, 453 U.S.\ 654 (1981):} Relying on \textit{Curtiss-Wright} and \textit{Youngstown}, the Court upheld executive orders suspending U.S.\ nationals' claims against Iran to resolve the hostage crisis, emphasizing longstanding congressional acquiescence in executive claims settlement while stressing the narrowness of its holding.
  \item \textbf{\textit{Zivotofsky v. Kerry}, 576 U.S.\ 1 (2015)} (\href{https://supreme.justia.com/cases/federal/us/576/1/}{Justia}): Although the Court held (6--3) that the President possesses the exclusive power to recognize foreign sovereigns---invalidating a statute requiring ``Israel'' to be listed on passports of citizens born in Jerusalem---Justice Kennedy's majority opinion \textbf{expressly repudiated} \textit{Curtiss-Wright}'s broad ``sole organ'' dicta. The Court explained that \textit{Curtiss-Wright} addressed a congressional delegation and declined to endorse an unbounded, extraconstitutional presidential monopoly over foreign policy, grounding the recognition power instead in Article~II's Reception Clause (\S\,3) and functional unity.
  \item \textbf{\textit{Trump v. Hawaii}, 585 U.S.\ 667 (2018):} Illustrated \textit{Curtiss-Wright}'s enduring vitality in statutory delegation and national security deference, upholding presidential entry restrictions under the broad text of \S\,212(f) of the Immigration and Nationality Act.
  \item \textbf{\textit{FCC v. Consumers' Research}, 606 U.S.\ 656 (2025)} (\href{https://supreme.justia.com/cases/federal/us/606/24-354/}{Justia}): In his concurring opinion, Justice Kavanaugh cited \textit{Curtiss-Wright} (299 U.S.\ at 319--322) to underscore that both the nondelegation doctrine and the major questions canon operate with sharply reduced force in national security and foreign policy domains because of the President's independent Article~II responsibilities.
  \item \textbf{\textit{Learning Resources, Inc. v. Trump}, 607 U.S.\ 229 (2026)} (\href{https://supreme.justia.com/cases/federal/us/607/24-1287/}{Justia}; \href{https://www.supremecourt.gov/opinions/25pdf/24-1287_4gcj.pdf}{Slip Op.}): \textit{Curtiss-Wright} took center stage in the Court's 6--3 invalidation of President Trump's tariffs imposed under IEEPA. Chief Justice Roberts's plurality (joined by Gorsuch and Barrett, JJ.) rejected the Government's and Justice Kavanaugh's proposed ``foreign affairs exception'' to the major questions doctrine, holding that peacetime tariffs belong to ``Congress alone'' under Article~I, \S\,8, cl.\ 1, rather than to concurrent or independent Article~II authority. Justice Gorsuch's concurrence explicitly cabined \textit{Curtiss-Wright} to its narrow statutory context---a targeted arms embargo on two warring nations conditioned on consultation and factual findings---refusing to read it as allowing Congress to hand off its core taxing and foreign-commerce powers through vague text. By contrast, Justice Kavanaugh's principal dissent (joined by Thomas and Alito, JJ.) and Justice Thomas's separate dissent leaned heavily on \textit{Curtiss-Wright} (299 U.S.\ at 315, 319, 321--322, 324) to argue that foreign-commerce and tariff delegations lie outside domestic major-questions skepticism and nondelegation constraints.
\end{itemize}

\section{Pedagogical Critiques and Exam Hypotheticals}
\textbf{Pedagogical Critiques:}
\begin{itemize}
  \item \textit{Flawed Historical Narrative:} Constitutional historians (notably Charles Lofgren, \textit{83 Yale L.J.\ 1 (1973)}, and David Levitan, \textit{55 Yale L.J.\ 467 (1946)}) have thoroughly critiqued Sutherland's ``extraconstitutional sovereignty'' thesis. Between 1776 and 1789, individual States frequently exercised external attributes (maintaining navies, imposing tariffs, and negotiating border agreements), and the Constitution itself carefully \textit{divides} foreign-affairs powers between Congress (declaring war, regulating foreign commerce, raising armies, funding diplomacy) and the President (Commander in Chief, negotiating treaties with Senate consent, receiving ambassadors).
  \item \textit{Misquoting John Marshall's ``Sole Organ'' Speech:} When Representative Marshall called the President the ``sole organ of the nation in its external relations'' in 1800, he was defending President Adams's execution of an existing extradition treaty---describing the President as the sole \textit{communicator} of national policy already established by law or treaty, not as the sole \textit{maker} of foreign policy. Sutherland transformed Marshall's modest conduit metaphor into a sweeping claim of substantive executive exclusivity.
  \item \textit{Logical Leap from Federal Power to Presidential Power:} Even if external sovereignty passed directly from the Crown to the United States as a whole, that historical premise explains why \textit{the federal government} (as against the States) holds plenary foreign-affairs power---it does not explain why that power belongs to \textit{the President} rather than \textit{Congress}.
\end{itemize}

\textbf{``Opposite Side'' Contrast (\textit{Curtiss-Wright} vs.\ \textit{Schechter Poultry} and \textit{Learning Resources}):}
Compare \textit{Curtiss-Wright} (1936) with \textit{Schechter Poultry} (1935) and \textit{Learning Resources, Inc. v. Trump} (2026). In \textit{Curtiss-Wright}, Justice Sutherland---a member of the conservative ``Four Horsemen'' who had just struck down domestic New Deal delegations in \textit{Panama Refining} and \textit{Schechter}---pivoted to strong executive deference because the statute regulated arms sales to foreign belligerents and implicated diplomatic confidentiality. Yet \textit{Learning Resources} (2026) demonstrates the modern boundary of that contrast: when an executive action touches foreign affairs \textit{only instrumentally} while exercising a core domestic-facing Article~I power vested exclusively in Congress (such as levying peacetime tariffs on domestic importers), a majority of the Court refuses to extend \textit{Curtiss-Wright}'s foreign-affairs carve-out to bypass clear congressional authorization.

\textbf{Exam Hypothetical:}
Suppose Congress passes the ``Global Semiconductor Security Act,'' authorizing the President, upon declaring a ``geopolitical supply-chain threat,'' to ``prohibit, license, or impose financial conditions on the export of advanced chips to any foreign nation if the President determines such action will advance U.S.\ strategic stability.'' Invoking the Act, the President issues a proclamation banning chip exports to Country~X and imposing a 20\% export fee on chips shipped to Country~Y, without consulting Congress or holding hearings. A U.S.\ chipmaker indicted for exporting to Country~X and assessed the fee on shipments to Country~Y challenges both measures.
\begin{itemize}
  \item \textit{Analysis of the Export Ban (Country~X):} Strong \textit{Curtiss-Wright} Zone~1 territory. Like the 1934 Chaco arms embargo, the statute targets a defined class of strategic goods in external relations, requires a presidential finding tied to foreign strategic stability, and implicates confidential national-security assessments. Under \textit{Curtiss-Wright} and Justice Kavanaugh's concurrence in \textit{Consumers' Research} (2025), a nondelegation challenge to the export ban will almost certainly fail.
  \item \textit{Analysis of the 20\% Export Fee (Country~Y):} Dramatically weaker for the President. First, even under \textit{Curtiss-Wright}, Sutherland expressly stated that executive foreign-affairs power ``must be exercised in subordination to the applicable provisions of the Constitution'' (299 U.S.\ at 320)---and the Export Clause (Art.\ I, \S\,9, cl.\ 5) flatly forbids taxes or duties on exports (a constitutional defect Chief Justice Roberts highlighted in \textit{Learning Resources} (2026)). Second, even if characterized as a regulatory condition rather than a tax, the chipmaker will invoke \textit{Learning Resources} (2026) and \textit{Zivotofsky} (2015) to argue that ``financial conditions'' does not clearly authorize revenue-raising exactions and that \textit{Curtiss-Wright} cannot override explicit Article~I limits.
\end{itemize}

\end{document}
